
\chapter{Мультивалютность и трансграничные платежи}\label{ch:multicurrency}
Один трансграничный платёж проходит через несколько
валют: покупатель видит цену в~одной, платит картой
или счётом в~другой, продавец получает выплаты в~третьей,
а~расчёт с~эквайером приходит в~четвёртой. Каждая сумма
фиксируется в~своём моменте времени. Для российского читателя
мультивалютность 2022--2026~годов -- работа под санкциями:
ряд банков отключён от SWIFT, корреспондентские
отношения перестраиваются вокруг китайских, белорусских,
казахстанских банков-посредников, доля расчётов
в~юанях и~других валютах дружественных стран растёт. Корреспондентские
счета и трансграничная логика карточных сетей разобраны
в~предыдущих частях; здесь -- архитектура мультивалютного
продукта и валютное регулирование после 2022~года.

\section{Кто управляет курсом}\label{sec:mc-currencies}

\highlight{В~мультивалютном продукте курс кто-то фиксирует, и~эта
сторона потом отвечает перед поддержкой и казначейством}.

В~первом варианте курс остаётся на стороне эмитента и~правил
карточной сети. Так устроен стандартный карточный сценарий,
когда продавец выставляет цену в~валюте операции, а~сумма
в~валюте карты определяется на стороне эмитента по правилам
валюты списания и конверсии~\cite{visa-core-rules,mc-dcc-guide}.
Для продавца это просто, но ни продавец, ни платформа
заранее не знают итоговую сумму списания в~валюте карты.

Во втором варианте курс контролирует эквайер или провайдер DCC,
который показывает держателю карты итоговую сумму в~домашней
валюте ещё до авторизации. Цена становится предсказуемой
для покупателя, но право задать курс и~удержать маржу
переходит к~стороне, которая управляет DCC~\cite{visa-core-rules,mc-dcc-guide}.

В~третьем варианте курс задаёт сама платформа. Так устроен
кошелёк, сервис путешествий или маркетплейс с~собственным
валютным модулем, где интерфейс показывает курс, платформа берёт
на себя риск его устаревания на коротком окне, а~затем
хеджирует или неттирует фактические потоки
на этапе расчёта.

В~четвёртом варианте курс приходит от внешнего казначейского
или валютного провайдера, а платёжный продукт привязывает
платёж к~уже рассчитанной сделке. Сценарий типичен
для B2B-сервисов и~платформ выплат, где цена перевода живёт
дольше одной сессии оплаты.

Проблема устаревшего курса возникает так: платформа
показывает клиенту курс в~момент выбора, подтверждение
приходит через час, и~за это время курс уже сдвинулся. Разница
между показанным и фактическим курсом -- убыток той стороны,
которая зафиксировала обязательство. Обычное решение состоит в~том,
что курс действует в~пределах ограниченного TTL, и при позднем
подтверждении платформа либо пересчитывает по актуальному курсу,
либо принимает риск на себя. Значения TTL и допустимый диапазон
волатильности для основных валютных пар оператор подбирает
по собственным данным; короткий TTL увеличивает долю
истёкших сессий, длинный создаёт валютный риск. Без явного
TTL потери на устаревшем курсе обнаруживаются только
при ежемесячной сверке казначейства.

\begin{figure}[tbp]
  \centering
  \includefiguresvg[width=.8\linewidth]{assets/figures/ch29-multicurrency/fx-ownership-models}
  \caption{Четыре модели владения курсом.}\label{fig:fx-ownership-models}
\end{figure}

Подписка с~ежемесячным списанием в EUR при выплате продавцу в USD
проходит через конверсию 12 раз в~год, и при каждой конверсии
курс другой; совокупная дельта составляет заметную
долю маржи платформы. Поэтому курс хранят как
отдельный объект с~явным источником, значением, моментом
фиксации, TTL и~стороной, несущей риск.

Минимальный реестр котировок принимает котировки от разных
источников, отбрасывает устаревшие (по TTL) и~слишком широкие
по спреду (выше заданного порога в~bps), и~возвращает самую
свежую из оставшихся. Источник времени вынесен в~параметр
\texttt{now}, чтобы тестируемое решение не зависело от системного
таймера:

\codefile{Python}{samples/ch34-fx-quote/quote.py}

\practicenote{%
  Если продукт обещает клиенту \enquote{ваш курс}, у~этого курса должны
  быть явный источник и~окно действия.

}
\section{Где фиксируется валюта расчёта}\label{sec:mc-settlement}

\textbf{Валюта цены (transaction currency)} -- та,
которую видит покупатель и которую продавец ожидает
на странице оплаты.
\textbf{Валюта списания (billing currency)} -- та,
в~которой эмитент дебетует держателя карты.
\textbf{Валюта расчёта (settlement currency)} -- та,
в~которой сеть или эквайер фактически рассчитывается с~вашей
стороной~\cite{visa-core-rules,mc-rules}.

В~цепочке участвуют три стороны с~разными валютами.
Держатель карты работает в~своей локальной валюте, расчётная
инфраструктура схемы -- в~валюте расчёта, эквайер --
в~своей локальной валюте. Исторически в~Visa и~Mastercard
валютой расчёта были USD или EUR; в~платёжной системе
\enquote{Мир} внутрироссийский расчёт идёт в~рублях через
платёжную систему Банка России. Во внутренней операции все три
валюты совпадают; в~трансграничной их расхождение неизбежно.

При \textbf{единой валюте расчёта} внутренний реестр обязательств
и выплаты ведутся в~одной базовой валюте, поэтому операционная
модель упрощается до одного казначейского баланса, одной
логики комиссий и меньшего числа сценариев для сверки. Цена
простоты -- постоянная конверсия на входе или выходе из системы
и зависимость от валютной наценки внешнего провайдера.

При \textbf{нескольких валютах расчёта} система хранит обязательства
и расчёты в~нескольких валютах. Реестр обязательств и логика
выплат сложнее, но лишней конверсии нет, а~реальные валютные
позиции не смешиваются в~одну синтетическую сумму.

Вопрос при проектировании: \enquote{EUR здесь --
  валюта цены, валюта списания,
валюта расчёта, валюта хранения обязательства или всё сразу?}
Ответы для разных слоёв системы часто не совпадают.

Продавец из России продаёт цифровой продукт покупателю
из Казахстана; цена выставлена в~рублях, покупатель
оплачивает картой казахстанского банка в~тенге, эквайер
рассчитывается с~продавцом в~рублях через расчётный
центр в~России. Валюта цены -- RUB, валюта списания -- KZT (валюта
карты), валюта расчёта -- RUB. Поскольку Visa и Mastercard
с 2022~года не работают для российских банков по новым
трансграничным операциям, маршрут авторизации казахстанской карты
к~продавцу в~России обычно идёт через UnionPay
(раздел~\ref{sec:up-international}), сети белорусской инфраструктуры
или соединения локальных платёжных систем
по двусторонним соглашениям. Реестр
обязательств должен хранить все точки фиксации -- сумму
в KZT на момент авторизации, применённый курс при клиринге
и~сумму в RUB на момент расчёта.

\begin{figure}[tbp]
  \centering
  \includefiguresvg[width=.8\linewidth]{assets/figures/ch29-multicurrency/three-currencies-timeline}
  \caption{Три валюты одной операции: цена, расчёт, выплата.}\label{fig:three-currencies-timeline}
\end{figure}

\section{DCC и маршрутизация по коридорам}\label{sec:mc-dcc}

DCC показывает держателю карты сумму в~знакомой валюте.
Расчёт по сети при этом идёт по своему маршруту, а~сумма
списания для клиента считается отдельно~\cite{visa-core-rules,mc-tpr}.

Продукт хранит исходную сумму, показанную сумму и применённый
курс как отдельные значения, не выводя одно из другого
задним числом. Без этого нельзя объяснить клиенту
списание и~сверить удержанную маржу DCC~\cite{visa-core-rules,mc-tpr}.

DCC зарабатывает на спреде, когда конверсия предлагается по курсу
хуже межбанковского, а~разница остаётся у~эквайера или его
провайдера DCC. Конкретное распределение определяется
коммерческим договором между ними и~в~публичных правилах схем
не закреплено. Держатель видит сумму в~знакомой
валюте, но платит за это удобство наценкой, встроенной
в~курс. Схемы требуют явного раскрытия наценки и~согласия
клиента до операции.

Иллюстрация механики DCC. Покупатель платит 100~EUR
картой европейского банка в~ТСП в~России. Без DCC эквайер
отправляет транзакцию на 100~EUR; сеть
конвертирует её при клиринге по своему оптовому курсу
с~небольшой наценкой, и~держатель видит итоговую сумму
в~своей домашней валюте только после списания. С DCC эквайер
через провайдера DCC предлагает держателю фиксированную
сумму в~его домашней валюте сразу на терминале, по курсу
с~заметно большей наценкой; разница между этой наценкой
и~оптовым курсом -- и~есть цена \enquote{удобства} видеть знакомую
валюту. Значения наценок отличаются у~разных
провайдеров DCC и публикуются в~раскрытии перед операцией.

Дерево решений для мультивалютного ТСП. Если ТСП продаёт
только в~одной стране и~получает выплаты в~местной
валюте -- DCC не нужен, конверсию делает сеть при клиринге
(самый дешёвый путь). Если ТСП хочет предложить DCC ради
повышения маржи -- нужен провайдер DCC, раскрытие наценки
и~согласие клиента. У~этого варианта высокий репутационный риск. Если ТСП
работает в~нескольких странах и~хочет минимизировать
трансграничные отказы -- локальный эквайер в~каждой стране
даёт максимальную долю одобрений и убирает трансграничную
наценку, но требует юридического присутствия и отдельной
расчётной модели в~каждой юрисдикции~\cite{visa-core-rules,mc-tpr}.

DCC -- часть политики маршрутизации. По коридору, BIN,
стране эквайера и~типу ТСП определяется, допустим ли DCC
и~откуда брать курс.

Согласие клиента, раскрытая наценка и~след аудита фиксируются
как поля события транзакции, иначе позже трудно доказать,
что конверсия была предложена корректно~\cite{visa-core-rules,mc-tpr}.

Платформа, работающая с~несколькими странами, эквайерами
и~схемами выплат, сталкивается с~задачей маршрутизации
по коридорам: для каждой пары \enquote{страна покупателя
/ страна ТСП} нужно выбрать эквайера, валютную схему
и~способ выплаты.

Выбор коридора зависит от стоимости авторизации и расчёта,
доли одобрений для нужного диапазона BIN, валюты выплаты
продавцу и~стороны, которая держит валютный риск.

В~базовом компромиссе дорогой маршрут через \enquote{родной}
эквайер держателя повышает \(P(\text{approve})\), но
увеличивает комиссию. Дешёвый маршрут через локального
эквайера снижает комиссию, но снижает \(P(\text{approve})\)
на трансграничных картах. При высоком среднем чеке потеря
на отказанной транзакции дороже разницы в~комиссии,
и~выбор смещается к~дорогому маршруту; при
низком чеке выгоднее обратное.

При \textbf{выплате в~домашней валюте продавца} платформа сама
консолидирует и~конвертирует все входящие потоки. Продавцу
проще, казначейской функции платформы -- сложнее.

При \textbf{выплате в~валюте исходного потока} платформа хранит
отдельные валютные карманы и~платит продавцу в~той валюте,
в~которой пришёл поток. Конверсий меньше, но продавец живёт
с~несколькими валютными балансами и отдельным графиком
выплат по каждому.

\begin{figure}[tbp]
  \centering
  \includefiguresvg[width=.75\linewidth]{assets/figures/ch29-multicurrency/corridor-routing}
  \caption{Четыре решения в~маршрутизации по коридорам.}\label{fig:corridor-routing}
\end{figure}

\practicenote{%
  Если внутри принимаются потоки в~нескольких валютах, а~продавцу обещан один
  баланс, у~валютной позиции, лимитов и источников курса должен
  быть явный владелец.

}
\section{Сверка по валютам}\label{sec:mc-fx-risk}

Расхождения в~мультивалютном продукте появляются после
авторизации, когда одна и~та же операция представлена сразу
в~нескольких видах: на странице оплаты -- в~валюте цены,
в~журнале авторизации -- одной суммой,
в~клиринге -- уже с~другой валютой расчёта, в~отчёте
о~выплате -- после удержания комиссии и~валютного спреда.
Валютная сверка остаётся отдельным процессом, и карточные правила
прямо требуют сверять данные сети с~внутренними записями,
а в Mastercard это сформулировано как ежедневная обязанность
сверять итоги и транзакции с~внутренним учётом~\cite{mc-rules}.

\highlight{Первое правило~--- сохранять исходную валютную пару. Внутренний
реестр обязательств хранит исходную сумму, сумму после
конверсии, источник курса и~момент его фиксации}.

Второе правило~--- считать комиссию отдельно от конверсии.
Если курсовая разница смешана с~комиссией эквайера,
понять, куда делась маржа, операционная команда не сможет.

Третье правило -- сверять выплаты по валютным карманам, а~не
по агрегированной сумме. Сходящаяся общая цифра скрывает
недостачу в~отдельной валюте, когда одна валюта компенсирует
другую в~общем итоге.

Пример валютной сверки одной операции выглядит так. Покупка за 100~EUR,
расчёт по курсу 1,08 = 108~USD, комиссия эквайера
2,5\% = 2,70~USD, interchange 1,5\% = 1,62~USD, комиссия
платформы 3\% = 3,24~USD, выплата продавцу =
108 $-$ 2,70 $-$ 1,62 $-$ 3,24 = 100,44~USD. Если продавец
получает выплату в GBP по курсу 0,79, итого 79,35~GBP. Одна
и~та же покупка существует в~пяти представлениях -- 100~EUR на
странице оплаты, 100~EUR в~авторизации, 108~USD в~расчёте,
100,44~USD в~нетто-выплате, 79,35~GBP на счёте продавца.

Ошибки сверки чаще всего возникают при смешении валютной
дельты и~комиссий. Если 7,70~USD удержаний
(2,70 + 1,62 + 3,24) учтены как курсовой убыток вместо
раздельных комиссий, казначейство видит завышенную курсовую
дельту и~ищет проблему там, где её нет. Каждое удержание
имеет свой тип: interchange, assessment, комиссия платформы,
валютная наценка~--- и~сверяется по типу; итоговая сумма для
этой сверки не подходит.

\begin{figure}[tbp]
  \centering
  \includefiguresvg[width=.85\linewidth]{assets/figures/ch29-multicurrency/multicurrency-recon}
  \caption{Одна транзакция в~пяти валютных представлениях.}\label{fig:multicurrency-recon}
\end{figure}

\section{Валютное регулирование после 2022~года}\label{sec:mc-russia}

Базовое регуляторное поле валютных операций в~России --
федеральный закон 173-ФЗ \enquote{О~валютном регулировании
и валютном контроле}~\cite{fz-173}. Он определяет понятия
резидента и нерезидента, перечень разрешённых валютных
операций между ними и~порядок валютного контроля. Подзаконные
акты Банка России конкретизируют процедуры представления
документов валютного контроля, постановки контрактов
на учёт в~уполномоченных банках и репатриации валютной
выручки~\cite{cbr-181i}.\footnote{Точные реквизиты
действующих Указаний и Инструкций Банка России (например,
Инструкция Банка России от 16.08.2017 № 181-И, заменившая ранее действовавшую Инструкцию 138-И)
сверяются по сайту регулятора перед публикацией, поскольку
порядок валютного контроля в~РФ менялся неоднократно после 2022~года.}
Для платёжного продукта, который проводит трансграничные
операции с~участием российской стороны, ссылки на 173-ФЗ
и подзаконные акты появляются в~нескольких точках -- при
открытии валютных счетов, при идентификации клиента, при
направлении валютной операции в~банк, при формировании
отчётности.

После марта 2022~года конфигурация международных валютных
расчётов для российских банков изменилась, когда ряд
крупнейших банков был отключён от SWIFT по решению Совета ЕС,
оформленному Регламентом Совета (EU) 2022/345~\cite{eu-swift-russia-2022,ofac-russia-sanctions}.
SWIFT работает для значительной части российских банков
и~для большинства корреспондентских отношений с~банками
дружественных стран, но для отключённых банков критическую
часть инфраструктуры пришлось перенести на альтернативные каналы.

\textbf{СПФС} (см.~\ref{sec:banks-infrastructure}) для отключённых
от SWIFT банков превратилась из локальной альтернативы в~основной
канал финансовых сообщений с~банками-партнёрами в~ЕАЭС, СНГ
и БРИКС~\cite{cbr-spfs}.

\textbf{Расчёты в~юанях} стали ключевым инструментом
для экспорта и~импорта после 2022~года.
Технически расчёты в CNY между российскими и китайскими
банками идут через систему CIPS (Cross-Border Interbank
Payment System) -- китайскую инфраструктуру межбанковских
расчётов в~юанях. Оператор -- Cross-border Interbank Payment
System Co., Ltd., под надзором Народного банка
Китая~\cite{cips-overview}. Российские банки
выходят в CIPS либо напрямую (как прямые или непрямые участники),
либо через банки-посредники в~Китае.

\textbf{Корсчета через банки-посредники.} Для отключённых
от SWIFT российских банков корреспондентские отношения
с~зарубежными банками-партнёрами перестраиваются через
банки-посредники в~дружественных юрисдикциях -- прежде всего
Китай, Индия, ОАЭ, Турция, страны Центральной Азии. Цепочка
расчётов удлиняется, каждая операция получает дополнительные
шаги по комплаенсу и валютному контролю.

\textbf{BRICS Pay и связанные инициативы.} На уровне
БРИКС обсуждаются проекты межоператорной интеграции
национальных платёжных систем (BRICS Pay, BRICS Bridge)
для прямых расчётов в~национальных валютах без обращения
к~доллару~\cite{brics-pay-overview}.\footnote{Текущий статус,
состав участников и техническая зрелость этих инициатив
сверяются по официальным сообщениям БРИКС и~Банка России перед
публикацией. Отдельный проект межбанковского обмена
CBDC mBridge ведётся под BIS Innovation Hub совместно
с~центральными банками Китая, ОАЭ, Гонконга и~Таиланда
(BIS вышел из проекта 31 октября 2024 года) и формально
к~повестке БРИКС не относится.}

Для архитектора платёжного продукта в~России отсюда
следуют три практических последствия. Во-первых, мультивалютная
модель должна предусматривать не одну инфраструктуру
SWIFT, а несколько параллельных каналов -- SWIFT (для
неотключённых банков и зарубежных коридоров), СПФС
(для внутреннего периметра), CIPS (для расчётов в~юанях),
двусторонние и многосторонние соглашения. Во-вторых,
выбор валют расчёта смещён к~национальным валютам
дружественных стран (CNY, AED, INR, KZT, BYN); USD и EUR
остаются возможными, но требуют согласования с~банком-партнёром
на каждой операции и подпадают под санкционный контроль.
В-третьих, валютный контроль (173-ФЗ) требует представления
документов и постановки контрактов на учёт независимо
от выбранного канала расчёта.

\practicenote{%
  Если продукт обещает клиенту \enquote{трансграничный
  платёж в~любой валюте}, такое обещание
  редко выполнимо без оговорок, поскольку набор реально работающих
  валют, банков-посредников и~каналов расчётов меняется
  и~зависит от санкционного статуса участников. Маркетинг
  должен идти за реальностью валютных коридоров; обратный порядок
  быстро рассыпается.

}

\section{Корреспондентский банкинг и ISO~20022}\label{sec:mc-correspondent}

Корреспондентские расчёты разобраны в~разделе~\ref{sec:banks-correspondent},
миграция к ISO~20022 -- в~разделе~\ref{sec:iso20022}.

Первый тезис -- расчёт по трансграничной операции почти всегда
проходит по более длинной цепочке, чем авторизация, и~идёт
в~другом темпе~\cite{fatf-correspondent-banking,swift-instant-interoperability}.
Одобрение за секунду не означает выплату в~тот же день.

Второй тезис -- расчётный банк, карточная сеть и~валютный
провайдер бывают тремя разными сторонами, поэтому продукт должен
явно фиксировать, кто отвечает за каждую стадию, без отговорок
вида \enquote{деньги идут через SWIFT}.

Третий тезис -- переход отрасли на ISO~20022 делает сообщения
богаче, но не снимает продуктовых решений. Структурированные
поля облегчают комплаенс-проверки и~сверку, но сами по себе
не отвечают, кто несёт валютный риск и~в~какой валюте
учитываются внутренние обязательства~\cite{iso-20022,swift-iso20022-milestone,swift-transaction-manager-iso20022}.
